\documentclass{article}
\usepackage[a4paper, total={16cm, 26cm}]{geometry}
\usepackage{texmaths-main}
\usepackage{enumitem}

\begin{document}
\begin{center}
    \Large Évaluation de cours : Statistiques/Équations (Sujet A)\par
	{\itshape\normalsize Répondre directement sur le sujet. Les réponses doivent être
	justifiées par un calcul. La calculatrice est autorisée.}
\end{center}
\vskip10pt
Nom : \hskip50pt Prénom : \hskip50pt Classe~:
\vskip20pt
% PARTIE I : Statistiques
\begin{enumerate}[label=\arabic*), font=\bfseries]
	\item Voici une série de nombres~:$$6,5\qquad13\qquad16,5\qquad18\qquad6$$
	\begin{enumerate}[label=\alph*), font=\bfseries]
		\item Calculer la moyenne.\dotedline\dotedline\dotedline
		\item Calculer la médiane.\dotedline\dotedline
		\item Calculer l'étendue.\dotedline
	\end{enumerate}
	\item Voici une série statistique~:
	\begin{center}
		\begin{tblr}{colspec=*6{Q[c,m]},vlines,hlines}
			Valeures&16,5&12,1&10,6&11,9&7,5\\
			Effectif&7&9&3&3&8
		\end{tblr}
	\end{center}
	\begin{enumerate}[label=\alph*), font=\bfseries]
		\item Calculer l'effectif total.\dotedline
		\item Calculer la moyenne.\dotedline\dotedline
		\item Calculer la médiane\dotedline\dotedline
		\item Calculer l'étendue.\dotedline
	\end{enumerate}
	\item Résoudre les problèmes suivants.
	\begin{enumerate}[label=\alph*), font=\bfseries]
		\item Dans une entreprise de 120 salariés, 65\% des salariés perçoivent
		un salaire de plus de 2800 \euro.\par
		Quel est le nombre de salariés rémunérés plus de 2800 \euro~?%
		\dotedline\dotedline\dotedline\dotedline
		\item Dans un collège, 210 élèves sont externes. Sachant que le nombre total
		d'élèves est de 480, quel est le pourcentage d'élèves externes~? (arrondir à
		l'unité).\dotedline\dotedline\dotedline\dotedline
		\item Un pull coûtant 55\euro\ voit son prix augmenter de 20 \%. Quel est son
		nouveau prix~?\dotedline\dotedline\dotedline
	\end{enumerate}
	\item Résoudre en $x$ les équations suivantes (détailler les étapes de calcul)~:\par\noindent
	\hbox{\hbox to.3\linewidth{$5x=17,5$\hfill}\hbox to.3\linewidth{$9x+5=6$\hfill}\hbox{$\dfrac32x-5=3$\hfill}}\par\noindent\vfill
	\hbox{\hbox to.3\linewidth{$-2x=6$\hfill}\hbox to.3\linewidth{$15-2x=20$\hfill}\hbox{$\dfrac45x+\dfrac67=\dfrac58$\hfill}}\par\vfill
	\item Le nombre $y=-4$ est-il solution de l'équation suivante~:$$(y+5)^2=5y+y^2$$\vskip100pt
\end{enumerate}
\newpage
\begin{center}
    \Large Évaluation de cours : Statistiques/Équations (Sujet B)\par
	{\itshape\normalsize Répondre directement sur le sujet. Les réponses doivent être
	justifiées par un calcul. La calculatrice est autorisée.}
\end{center}
\vskip10pt
Nom : \hskip50pt Prénom : \hskip50pt Classe~:
\vskip20pt
% PARTIE I : Statistiques
\begin{enumerate}[label=\arabic*), font=\bfseries]
	\item Voici une série de nombres~:$$12,5\qquad15\qquad17,5\qquad16\qquad8$$
	\begin{enumerate}[label=\alph*), font=\bfseries]
		\item Calculer la moyenne.\dotedline\dotedline\dotedline
		\item Calculer la médiane.\dotedline\dotedline
		\item Calculer l'étendue.\dotedline
	\end{enumerate}
	\item Voici une série statistique~:
	\begin{center}
		\begin{tblr}{colspec=*6{Q[c,m]},vlines,hlines}
			Valeures&18,5&12,3&8,6&12,9&11,5\\
			Effectif&3&7&5&8&2
		\end{tblr}
	\end{center}
	\begin{enumerate}[label=\alph*), font=\bfseries]
		\item Calculer l'effectif total.\dotedline
		\item Calculer la moyenne.\dotedline\dotedline
		\item Calculer la médiane\dotedline\dotedline
		\item Calculer l'étendue.\dotedline
	\end{enumerate}
	\item Résoudre les problèmes suivants.
	\begin{enumerate}[label=\alph*), font=\bfseries]
		\item Dans une entreprise de 120 salariés, 45\% des salariés perçoivent
		un salaire de plus de 2800 \euro.\par
		Quel est le nombre de salariés rémunérés plus de 2800 \euro~?%
		\dotedline\dotedline\dotedline\dotedline
		\item Dans un collège, 360 élèves sont externes. Sachant que le nombre total
		d'élèves est de 530, quel est le pourcentage d'élèves externes~? (arrondir à
		l'unité).\dotedline\dotedline\dotedline\dotedline
		\item Un pull coûtant 65 \euro\ voit son prix augmenter de 15 \%. Quel est son
		nouveau prix~?\dotedline\dotedline\dotedline
	\end{enumerate}
	\item Résoudre en $x$ les équations suivantes (détailler les étapes de calcul)~:\par\noindent
	\hbox{\hbox to.3\linewidth{$3x=28,5$\hfill}\hbox to.3\linewidth{$7x+3=3$\hfill}\hbox{$\dfrac72x-4=1$\hfill}}\par\noindent\vfill
	\hbox{\hbox to.3\linewidth{$-2x=6$\hfill}\hbox to.3\linewidth{$15-2x=20$\hfill}\hbox{$\dfrac85x+\dfrac34=\dfrac65$\hfill}}\par\vfill
	\item Le nombre $y=-3$ est-il solution de l'équation suivante~:$$(y+4)^2=6y+y^2$$\vskip100pt
\end{enumerate}
\end{document}